\chapter{Fossen's six-degree-of-freedom model}
\label{ch:fossen}

The starting point of \file{maneuvers.md} is the manoeuvring model of Fossen
\cite[ch.~6--7]{fossen2011}, written with the velocity relative to the water
$\vect{\nu}_r = \vect{\nu} - \vect{\nu}_c$:
\begin{equation}
  \underbrace{(\vect{M}_\mathrm{RB} + \vect{M}_\mathrm{A})\,\dot{\vect{\nu}}_r}_{\text{inertia}}
  + \underbrace{\vect{C}_\mathrm{RB}(\vect{\nu})\vect{\nu} + \vect{C}_\mathrm{A}(\vect{\nu}_r)\vect{\nu}_r}_{\text{Coriolis--centripetal}}
  + \underbrace{\vect{D}(\vect{\nu}_r)\,\vect{\nu}_r}_{\text{damping}}
  + \underbrace{\vect{g}(\vect{\eta})}_{\text{restoring}}
  = \vect{\tau}_\mathrm{thrust} + \vect{\tau}_\mathrm{dist}.
  \label{eq:fossen}
\end{equation}
This chapter takes the terms one at a time and explains what physical effect each
represents, how big it is, and which coefficient in \file{model.yaml} belongs to it.

\section{Rigid-body dynamics}

Newton's and Euler's laws about the centre of mass, written in a rotating body frame, read
\begin{align}
  m\,(\dot{\vect{v}} + \vect{\omega}\times\vect{v}) &= \vect{f}, \label{eq:newton}\\
  \vect{I}\,\dot{\vect{\omega}} + \vect{\omega}\times(\vect{I}\vect{\omega}) &= \vect{m}, \label{eq:euler}
\end{align}
with $\vect{v} = [u,v,w]^\top$ and $\vect{\omega} = [p,q,r]^\top$. The cross products are
not forces. They appear because the frame in which $\dot{\vect{v}}$ is measured rotates
with the vehicle. In matrix form, \cref{eq:newton,eq:euler} become
$\vect{M}_\mathrm{RB}\dot{\vect{\nu}} + \vect{C}_\mathrm{RB}(\vect{\nu})\vect{\nu}$. The surge
row of \eqref{eq:newton} is
\begin{equation}
  m\,(\dot u - v r + w q) = X .
  \label{eq:surge-rb}
\end{equation}
Both $v r$ and $w q$ are products of two components that flip together under the FRD/FLU
reflection, so \eqref{eq:surge-rb} reads identically in both conventions. The origin of
\code{base_link} lies at the centre of mass (per the DVL launch file), so the lever-arm terms
that Fossen carries for an offset centre of gravity vanish.

\section{Added mass}
\label{sec:addedmass}

When a body accelerates through a fluid, it must also accelerate some of the surrounding
fluid. The fluid's kinetic energy is a quadratic form of the body velocity,
$T_\mathrm{A} = \tfrac12\vect{\nu}^\top\vect{M}_\mathrm{A}\vect{\nu}$, and the energy needed to
build it up appears to the body as extra inertia \cite{newman1977}. For a sphere the added
mass is exactly half the displaced fluid mass. For an ROV-shaped body it is typically
\SIrange{30}{100}{\percent} of the dry mass in translation. In the BlueROV2 data,
$-\Xud = \qty{6.36}{\kilogram}$ against $m = \qty{13.5}{\kilogram}$.

The SNAME notation writes added mass as a \emph{hydrodynamic derivative}:
$\Xud = \partial X/\partial\dot u$, the change of hydrodynamic surge force per unit surge
acceleration. Accelerating forward causes the fluid to push \emph{back}, so $\Xud < 0$, and
moving it to the left-hand side gives
\begin{equation}
  (m - \Xud)\,\dot u = \dots, \qquad \mtot \coloneqq m - \Xud > m.
  \label{eq:mtot}
\end{equation}

\begin{caution}
Sign conventions differ between the code and much of the literature.
\begin{itemize}[nosep]
  \item \textbf{Added mass} is stored with Fossen's sign: \code{X_udot} is \emph{negative}.
        Tables that list a positive magnitude (as von Benzon et al.\ \cite{vonbenzon2022} do)
        must be negated when entered.
  \item \textbf{Damping} is stored as a positive \emph{magnitude}: \code{X_u} $> 0$, and the
        force is $-X_u u$. Fossen writes the force as $+X_u u$ with $X_u < 0$.
\end{itemize}
\code{SurgeParams.__post_init__} in \file{dynamics.py} rejects $\mtot \le 0$. That catches a
wrongly signed added mass that is large enough to make the inertia non-positive. A smaller
sign error does not violate the check: it gives a valid but wrong $\mtot < m$.
\end{caution}

\subsection{Added-mass Coriolis terms}

The added mass also contributes Coriolis-like terms. The fluid momentum
$\vect{p}_\mathrm{A} = \vect{M}_\mathrm{A}\vect{\nu}$ is carried along by the rotating frame
in the same way as the body momentum (Kirchhoff's equations \cite{newman1977}). For a diagonal
$\vect{M}_\mathrm{A} = -\mathrm{diag}(\Xud, Y_{\dot v}, Z_{\dot w}, \dots)$ the translational
part of the fluid's reaction is $-\vect{\omega}\times\vect{p}_\mathrm{A}$, whose $x$
component is
\begin{equation}
  -\bigl(q\,(-Z_{\dot w}w) - r\,(-Y_{\dot v}v)\bigr) = Z_{\dot w}\,wq - Y_{\dot v}\,vr .
\end{equation}
Adding this to \eqref{eq:surge-rb} gives the classical surge equation with its two
velocity-product terms:
\begin{equation}
  \boxed{\;(m - \Xud)\,\dot u \;-\; (m - Y_{\dot v})\,v r \;+\; (m - Z_{\dot w})\,w q \;=\; X_\mathrm{hydro} + X_\mathrm{thrust}\;}
  \label{eq:surge-full}
\end{equation}
With the BlueROV2 numbers $m - Y_{\dot v} = 13.5 + 7.12 = \qty{20.6}{\kilogram}$. A vehicle
that turns at $r = \qty{0.5}{\radian\per\second}$ while side-slipping at
$v = \qty{0.1}{\meter\per\second}$ therefore experiences about \qty{1}{\newton} of apparent
surge force. That is small but not negligible against a thrust of a few newtons. These are
exactly the \code{v*r} and \code{w*q} candidates in the coupling check (\cref{ch:coupling}),
and the synthetic test in \file{tests/test_coupling_check.py} injects $k\,vr$ with
$k = -20$ for this reason.

\section{Damping}

Hydrodynamic damping has several physical sources: skin friction, pressure (form) drag from
flow separation, vortex shedding, and lift on appendages. For manoeuvring at low speed it is
modelled as the sum of a linear and a quadratic part \cite[sec.~6.4]{fossen2011}:
\begin{equation}
  \vect{D}(\vect{\nu})\vect{\nu} = \vect{D}_l\,\vect{\nu} + \vect{D}_n(\vect{\nu})\,\vect{\nu},
  \qquad
  \text{surge: } \; X_\mathrm{hydro} = -X_u\,u - \Xuu\,|u|\,u .
  \label{eq:damping}
\end{equation}
\begin{itemize}
  \item The \emph{quadratic} term is the familiar drag law $F = \tfrac12\rho C_D A\,|u|u$
        at high Reynolds number. The form $|u|u$ rather than $u^2$ makes the force odd in
        $u$, so drag always opposes the motion.
  \item The \emph{linear} term dominates at very low speed (laminar skin friction, and
        effects such as wave radiation at the surface). Its absolute size is small. It
        matters mainly because it sets the decay rate near zero speed.
\end{itemize}
Off-diagonal damping terms are neglected by assumption. Their absence is part of what the
coupling check tests.

\section{Restoring forces}
\label{sec:restoring}

Gravity $W = mg$ acts at the centre of gravity $\vect{r}_g$ and buoyancy
$B = \rho g\nabla$ at the centre of buoyancy $\vect{r}_b$. In the ENU world frame gravity
points along $-z$. Building the restoring vector from vectors rather than from Fossen's NED
table gives the correct signs automatically:
\begin{equation}
  \vect{f}_g = \vect{R}^\top\begin{bmatrix}0\\0\\-W\end{bmatrix}, \quad
  \vect{f}_b = \vect{R}^\top\begin{bmatrix}0\\0\\B\end{bmatrix}, \quad
  \vect{g}(\vect{\eta}) = -\begin{bmatrix} \vect{f}_g + \vect{f}_b \\
     \vect{r}_g\times\vect{f}_g + \vect{r}_b\times\vect{f}_b \end{bmatrix}.
  \label{eq:restoring}
\end{equation}

\paragraph{Small-angle pitch stiffness.} Take $\vect{r}_g = \vect{0}$,
$\vect{r}_b = [0, 0, \overline{BG}]^\top$ (buoyancy \emph{above} gravity, so $z_b > 0$ in FLU,
whereas Fossen's NED tables have it negative) and a pure pitch $\vect{R} = \vect{R}_y(\theta)$.
Then $\vect{f}_b = [-B\sin\theta,\;0,\;B\cos\theta]^\top$ and
\begin{equation}
  \vect{r}_b\times\vect{f}_b = [\,0,\; -B\,\overline{BG}\sin\theta,\; 0\,]^\top
  \;\;\Rightarrow\;\;
  g_\theta = B\,\overline{BG}\sin\theta \approx M_\theta\,\theta,
  \quad M_\theta = B\,\overline{BG}.
  \label{eq:pitch-stiffness}
\end{equation}
The moment opposes the tilt, so the vehicle is statically stable. Roll gives the same result,
$K_\phi = B\,\overline{BG}$ (for $W = B$). The dummy file derives $K_\phi = M_\theta =
\qty{131.6}{\newton}\times\qty{0.01}{\meter} = \qty{1.316}{\newton\meter\per\radian}$ in
exactly this way.

\paragraph{Trim angles.} If the centre of buoyancy is not exactly above the centre of
gravity, the vehicle floats tilted at $\phi_0, \theta_0$, and the restoring torque pulls
towards that attitude rather than towards level: $g_\theta \approx M_\theta(\theta-\theta_0)$.
This is why \file{model.yaml} carries \code{phi_0} and \code{theta_0}, to be measured from
a still, powered-down float. (Since 15 October 2026 only \code{theta_0}; roll is not
modelled, \cref{sec:no-roll}.)

\paragraph{Yaw.} No restoring moment exists in yaw, because buoyancy and gravity are both
vertical. Yaw is a pure integrator with damping and never oscillates back. This is why the
static tests in \file{maneuvers.md} cover roll and pitch only.

\paragraph{Net buoyancy in surge.} If $W \ne B$, the translational part of
\eqref{eq:restoring} gives a surge component $(W - B)\sin\theta$ whenever the vehicle is
pitched. This is the \code{sin(pitch)} candidate of the coupling check.

\section{The model schema in \file{model.yaml}}
\label{sec:model-schema}

\begin{revised}[title={New section, added 15 October 2026}]
Until this revision \file{model.yaml} had slots for surge, roll, pitch and yaw only. The sway
and heave coefficients $Y_{\dot v}$, $Y_v$, $Y_{v|v|}$, $Z_{\dot w}$, $Z_w$, $Z_{w|w|}$ and the
buoyancy $B$ were added, so that the file now holds the complete diagonal model of
\eqref{eq:fossen}. No code reads the new entries yet.

\textbf{Later the same day} roll was removed from the model (\cref{sec:no-roll}): the slots
$I_{xx}$, $K_{\dot p}$, $K_p$, $K_{p|p|}$, $K_\phi$ and $\phi_0$ were deleted, the roll row of
\cref{tab:model-schema} was changed, and the equations below were written without roll.
\file{model.yaml} now has 21 parameters.
\end{revised}

With the port--starboard symmetry of \cref{ch:symmetry} and the off-diagonal terms neglected,
\eqref{eq:fossen} is fully determined by the parameters in \cref{tab:model-schema}. The
Coriolis--centripetal matrices $\vect{C}_\mathrm{RB}$ and $\vect{C}_\mathrm{A}$ need no
parameters of their own: they are built from $m$, the inertias and the same added-mass
derivatives \cite{fossen2011}. With the abbreviations
$m_{11} = m - \Xud$, $m_{22} = m - Y_{\dot v}$, $m_{33} = m - Z_{\dot w}$ and the roll
rate $p$ set to zero (\cref{sec:no-roll}), the translational equations are
\begin{align}
  m_{11}\,\dot u &= \tau_X - X_u\,u - \Xuu\,|u|\,u + (W - B)\sin\theta
                    + m_{22}\,v r - m_{33}\,w q, \label{eq:diag-surge}\\
  m_{22}\,\dot v &= \tau_Y - Y_v\,v - Y_{v|v|}\,|v|\,v + (B - W)\sin\phi\cos\theta
                    - m_{11}\,u r, \label{eq:diag-sway}\\
  m_{33}\,\dot w &= \tau_Z - Z_w\,w - Z_{w|w|}\,|w|\,w + (B - W)\cos\phi\cos\theta
                    + m_{11}\,u q. \label{eq:diag-heave}
\end{align}
The restoring terms follow from \eqref{eq:restoring} in the same way as
\eqref{eq:pitch-stiffness}: $\vect{R}^\top[0,0,1]^\top = [-\sin\theta,\ \sin\phi\cos\theta,\
\cos\phi\cos\theta]^\top$, and the net force on the vehicle is $(B-W)$ times this vector.
The roll angle $\phi$ is taken from the EKF; it is a measured input, not a state. Pitch and
yaw follow in the same way with $M_\theta(\theta-\theta_0)$ as the pitch restoring moment and
none in yaw. The Coriolis products in \eqref{eq:diag-surge}--\eqref{eq:diag-heave} are the
components of $-\vect{\omega}\times(\vect{M}_t\vect{v})$ with
$\vect{M}_t = \mathrm{diag}(m_{11}, m_{22}, m_{33})$. This form follows from vectors and holds
in any right-handed body frame, including FLU.

\begin{table}[htbp]
\centering
\caption{Parameters of the diagonal model in \file{model.yaml}, grouped by degree of freedom,
with the inertia, damping and restoring entries of each.}
\label{tab:model-schema}
\small
\begin{tabularx}{\textwidth}{@{}l l X l@{}}
\toprule
DOF & Inertia & Damping (\code{damping_linear}, \code{damping_quadratic}) & Restoring \\
\midrule
Surge & $m$, $\Xud$ & $X_u$, $\Xuu$ & via $W - B$ \\
Sway  & $m$, $Y_{\dot v}$ & $Y_v$, $Y_{v|v|}$ & via $W - B$ \\
Heave & $m$, $Z_{\dot w}$ & $Z_w$, $Z_{w|w|}$ & $B$ (with $W = mg$) \\
Roll  & \multicolumn{3}{l}{not modelled; $\phi$ and $p$ from the EKF (\cref{sec:no-roll})} \\
Pitch & $I_{yy}$, $M_{\dot q}$ & $M_q$, $M_{q|q|}$ & $M_\theta$, $\theta_0$ \\
Yaw   & $I_{zz}$, $N_{\dot r}$ & $N_r$, $N_{r|r|}$ & none \\
\bottomrule
\end{tabularx}
\end{table}
\FloatBarrier

\paragraph{What the schema deliberately leaves out.}
\begin{itemize}
  \item \textbf{Centre of gravity and centre of buoyancy.} They enter only through
        $M_\theta$ and $\theta_0$, which are lumped and measured directly.
        Keeping both $\vect{r}_b$ and $M_\theta$ would describe the same physics twice.
  \item \textbf{Off-diagonal terms} such as $X_{\dot w}$, $M_{\dot u}$ or $Z_q$. They are added
        only if the combined-axes runs leave a residual the diagonal model cannot explain
        (\cref{ch:symmetry}).
  \item \textbf{Water current and tether.} These are treated as disturbances, not model
        parameters.
\end{itemize}

\begin{caution}
$B - W$ is a small difference of two large forces. Weighing the dry vehicle ($m$) and computing
$B$ from the displaced volume will not give it to useful accuracy. A direct measurement of the
net buoyancy, for example with a spring scale on the submerged vehicle, is preferable. If
POLARIS operates at the surface, heave restoring is dominated by the waterplane stiffness
$\rho g A_\mathrm{wp}$, which has no slot yet.
\end{caution}

\subsection{Roll is not modelled}
\label{sec:no-roll}

\begin{revised}[title={New subsection, added 15 October 2026}]
The decision to drop roll and the evidence from the recorded bags. The roll statistics were
computed with a throw-away script over the EKF of the existing recordings; they are not
reproduced by \file{gen_figdata.py}.
\end{revised}

POLARIS is passively roll-stable: the centre of buoyancy lies above the centre of gravity,
and the restoring moment $K_\phi\phi$ of \eqref{eq:pitch-stiffness} pulls it back upright.
Roll is never commanded during identification. The roll equation is therefore dropped, and so
is every Coriolis product that contains the roll rate $p$. \Cref{tab:roll-terms} lists them.
As a consequence $I_{xx}$ and $K_{\dot p}$ no longer appear in any remaining equation, and
the slots $I_{xx}$, $K_{\dot p}$, $K_p$, $K_{p|p|}$, $K_\phi$, $\phi_0$ were removed from
\file{model.yaml}.

\begin{table}[htbp]
\centering
\caption{Coriolis--centripetal products containing the roll rate $p$ that are removed by not
modelling roll, from $-\vect{\omega}\times(\vect{M}_t\vect{v})$ and
$-\vect{\omega}\times(\vect{I}_t\vect{\omega}) - \vect{v}\times(\vect{M}_t\vect{v})$, with
$I_{11} = I_{xx} - K_{\dot p}$, $I_{22} = I_{yy} - M_{\dot q}$, $I_{33} = I_{zz} - N_{\dot r}$.}
\label{tab:roll-terms}
\small
\begin{tabular}{@{}lll@{}}
\toprule
Equation & Removed term & Physical meaning \\
\midrule
Sway  & $+m_{33}\,w p$ & Heave velocity turned into sway by rolling \\
Heave & $-m_{22}\,v p$ & Sway velocity turned into heave by rolling \\
Pitch & $-(I_{11} - I_{33})\,p r$ & Gyroscopic pitch moment of combined roll and yaw rates \\
Yaw   & $-(I_{22} - I_{11})\,p q$ & Gyroscopic yaw moment of combined roll and pitch rates \\
Roll  & whole equation & Roll inertia, damping and stiffness \\
\bottomrule
\end{tabular}
\end{table}
\FloatBarrier

\paragraph{Surge is unaffected.} The diagonal surge equation \eqref{eq:diag-surge} contains no
$p$ and no $\phi$. Dropping roll changes nothing in the surge work of \cref{ch:surge}.

\paragraph{The roll angle is still used.} Dropping the roll \emph{dynamics} is not the same as
setting $\phi = 0$. The angle enters the restoring terms of \eqref{eq:diag-sway} and
\eqref{eq:diag-heave} and the kinematics $\vect{J}(\vect{\eta})$ that turn body velocities into
world positions. The EKF provides $\phi$ at no cost, so it is replayed there, exactly as
\code{dofs.roll: measured} already did for the propagation (\cref{ch:symmetry}). Setting
$\phi = 0$ would discard information without simplifying anything.

\paragraph{Evidence from the recordings.} \Cref{tab:roll-evidence} summarises the EKF roll angle
of the existing bags, restricted to samples with $|u| > \qty{0.3}{\meter\per\second}$.
``Straight'' means $|r| < \qty{3}{\degree\per\second}$, ``turning'' means
$|r| > \qty{10}{\degree\per\second}$.

\begin{table}[htbp]
\centering
\caption{EKF roll angle in the existing recordings while moving: median roll, 95th percentile of
the deviation from the median on straight segments and in turns, and the dominant frequency of
the roll signal.}
\label{tab:roll-evidence}
\small
\begin{tabular}{@{}lcccc@{}}
\toprule
Bag & Median & $|\Delta\phi|_{95}$ straight & $|\Delta\phi|_{95}$ turning & Peak frequency \\
 & (\unit{\degree}) & (\unit{\degree}) & (\unit{\degree}) & (\unit{\hertz}) \\
\midrule
\file{zermatt_rectangle_01} & $-2.0$ & 2.0 & 3.9 & 0.56 \\
\file{zermatt_rectangle_04} & $-2.2$ & 2.1 & -- & 0.03 \\
\file{zermatt_rectangle_06} & $-2.1$ & 2.4 & -- & 0.03 \\
\file{zermatt_rectangle_07} & $-2.2$ & 2.0 & -- & 0.06 \\
\file{zermatt_rectangle_08} & $-2.0$ & 2.1 & 3.4 & 0.63 \\
\file{zermatt_grid_01}      & $-0.6$ & 3.1 & -- & 0.03 \\
\file{zermatt_grid_02}      & $-0.6$ & 6.7 & 88 & 0.03 \\
\file{aut_no_z_surface_big_40_grid_03} & $-1.1$ & 14.2 & 5.5 & 0.56 \\
\bottomrule
\end{tabular}
\end{table}
\FloatBarrier

Three observations follow from \cref{tab:roll-evidence}.
\begin{enumerate}
  \item In the Zermatt rectangle runs roll stays within about \qty{2}{\degree} on straight
        segments and about \qty{4}{\degree} in turns. Turn-induced heel is small, so the
        removed terms of \cref{tab:roll-terms} are small there as well.
  \item In the June surface run, roll reaches \qty{14}{\degree} on \emph{straight} segments, with
        a peak at \qty{0.56}{\hertz}. That is wave-induced roll, not a consequence of
        manoeuvring. A roll model would not predict it either, because the wave forcing is not
        measured. The roll rate there reaches about \qty{43}{\degree\per\second} (95th percentile), so the removed
        sway and heave products are no longer negligible. Such bags should not be used for sway
        or heave fits.
  \item The turning value of \file{zermatt_grid_02} (\qty{88}{\degree}) and the large maxima in
        some bags are almost certainly handling or recovery while the EKF was running, not
        manoeuvring. This was not checked further. In every bag the vehicle moves faster than
        \qty{0.3}{\meter\per\second} only \qtyrange{2}{42}{\percent} of the time.
\end{enumerate}
The consistent median of about \qty{-2}{\degree} in the rectangle runs is either a real list of
the vehicle or an IMU mounting offset. It does not matter for the model, because $\phi$ is
replayed from the EKF.

\begin{caution}
The decision is sound as long as (i) bags used for sway, heave or yaw fits have small roll, for
example $|\Delta\phi|_{95} < \qty{5}{\degree}$; (ii) ArduSub is not commanding roll during the
runs, since the flight mode may stabilise roll actively through the rear vertical pair; and
(iii) the combined surge--yaw runs of \cref{ch:symmetry} do not leave a residual correlated with
$p$. If any of these fails, the six slots can be restored from git history.
\end{caution}

\section{Thrust and disturbances}

$\vect{\tau}_\mathrm{thrust}$ is never measured directly. It is reconstructed from the
logged PWM through the thruster map (\cref{ch:actuation}) and the allocation matrix.
$\vect{\tau}_\mathrm{dist}$ collects everything that is not modelled: water current (which
enters through $\vect{\nu}_c$), tether drag, and waves when at the surface. For the
identification these are treated as \emph{nuisance} terms. A slowly varying current or
tether pull looks like a constant force offset, which is why \code{coupling_check.py} fits
an optional \code{offset [N]} regressor.

\begin{questions}
\begin{enumerate}
  \item Show from \eqref{eq:newton} that the surge equation contains $-m\,vr$ and $+m\,wq$.
  \item A colleague enters $\Xud = +\qty{6.36}{\kilogram}$ from a paper into \file{model.yaml}.
        What $\mtot$ results for $m = \qty{13.5}{\kilogram}$, and would
        \code{SurgeParams} reject it?
  \item Why does the restoring vector contain no yaw stiffness?
  \item Where in \file{model.yaml} would a measured $\overline{BG}$ end up?
\end{enumerate}
\end{questions}
